% Pandoc template for IEEE conference paper sections
\section{Introduction}

Pairing-friendly elliptic curves underpin a growing class of
cryptographic protocols, from BLS aggregate signatures \cite{BLS2001}
to succinct non-interactive zero-knowledge proofs (zk-SNARKs)
\cite{Groth2016, PLONK2019}. These protocols require curves whose
scalar field and base field support efficient polynomial arithmetic via
the Number Theoretic Transform (NTT), which in turn demands that both
field primes satisfy \(q \equiv 1 \pmod{2^s}\) for a sufficiently large
\(s\).

The Cocks-Pinch method \cite{CocksPinch2001} offers unmatched
flexibility in constructing pairing-friendly curves with arbitrary
embedding degrees. Unlike parametric families such as BN
\cite{FreemanScottTeske2010} or BLS12 \cite{BLS12_381}, the
Cocks-Pinch approach can target specific security levels and field
sizes. However, the traditional algorithm treats NTT-friendliness as an
afterthought---primes are generated with random two-adicity, typically
only 1--3 bits, rendering them unsuitable for modern proof systems.

A separate but related concern is \emph{parameter auditability}. The
Dual\_EC\_DRBG incident \cite{DualECDRBG} demonstrated that opaque
cryptographic constants can conceal backdoors. While the Cocks-Pinch
method is inherently transparent (all parameters derive from a public
seed), the resulting primes are typically inscrutable hexadecimal blobs
that resist manual inspection.

\textbf{Contributions.} This paper presents three results:

\begin{enumerate}
\def\labelenumi{\arabic{enumi}.}
\item
  \textbf{Improved Cocks-Pinch with NTT constraints.} We modify the
  Cocks-Pinch construction to \emph{guarantee} high two-adicity in both
  field primes by constraining the cyclotomic seed \(T\). Specifically,
  choosing \(T \equiv 0 \pmod{2^{12}}\) yields \(\nu_2(r-1) \geq 36\)
  deterministically (formalized in Lemma 1, Section 3), while an
  extended lift search ensures \(\nu_2(p-1) \geq 36\). In our
  construction, we obtain a two-adicity of 36, compared to BLS12-381's
  value of 32.
\item
  \textbf{Readability-aware prime selection.} We introduce a
  multi-criteria scoring system that evaluates candidate primes on eight
  structural properties (binary sparseness, zero limbs, hex density,
  etc.) and selects the most human-auditable parameters without
  sacrificing security.
\item
  \textbf{Curve1024: a concrete 1024-bit pairing-friendly curve.} We
  instantiate the above techniques to produce a complete, publicly
  verifiable curve with embedding degree \(k=18\), scalar field order
  \(r \approx 2^{512}\) (256-bit ECDLP security), and equation
  \(E: y^2 = x^3 + 41\). The full parameter set---including generator
  point---is published alongside a Rust implementation whose test suite
  verifies group arithmetic correctness and confirms resistance against
  MOV, Anomalous, and TNFS attacks.
\end{enumerate}

\textbf{Organization.} Section 2 reviews background on pairing-friendly
curves and the CM method. Section 3 presents our NTT-friendly
Cocks-Pinch modification. Section 4 describes the readability scoring
system. Section 5 presents the Curve1024 parameters and evaluates the
implementation. Section 6 discusses related work, and Section 7
concludes.
